\documentclass[11pt]{amsart}
\usepackage[T1]{fontenc}
\usepackage[margin=1.02in]{geometry}
\DeclareFontShape{OT1}{cmr}{bx}{n}{<-> cmbx10}{}
\DeclareMathSizes{8}{8}{8}{8}
\DeclareMathSizes{9}{9}{8}{8}
\DeclareMathSizes{10.95}{10.95}{8}{8}
\newtheorem{theorem}{Theorem}
\newcommand{\CIM}{\operatorname{CIM}}
\title{A five-vertex counterexample to the characteristic imset sink conjecture}
\date{November 27, 2026}
\begin{document}
\begin{abstract}
Making a vertex a sink can change the characteristic imset without
moving along a polytope edge. On a diamond with a pendant vertex,
four acyclic orientations give two distinct segments with the same
midpoint. This answers Conjecture~1 of Restadh in the negative.
Exact integer exposing functionals certify minimality on at most
four vertices. The example extends to connected graphs of every
order at least five. The earlier full five-vertex edge census
already classified the relevant pair.
\end{abstract}
\maketitle

\section{The conjecture}
For a directed acyclic graph \(H\) and vertex \(v\), write
\(\operatorname{pa}_H(v)\) for the vertices with arcs into \(v\).
The characteristic imset is the \(0/1\) vector whose coordinate
at \(S\), for each vertex subset \(S\) of size at least two, is
\[
 c_H(S)=1\quad\Longleftrightarrow\quad
 \exists v\in S:\ S\setminus\{v\}\subseteq\operatorname{pa}_H(v).
\]
In particular its pair coordinates record the skeleton. Two DAGs
have the same imset exactly when they are Markov equivalent
\cite[Theorem~1.4]{diam}. For an undirected \(G\), the convex hull
of the imsets of all DAGs with skeleton \(G\) is \(\CIM_G\).
Every distinct imset is a vertex. A convex combination of binary
vectors can equal a binary vector only when every positive-weight
summand agrees with it in every coordinate.

The sink operation \(H\downarrow i\) reverses every outgoing arc
at \(i\) and preserves all other arcs. It preserves acyclicity
because any new cycle would pass through the new sink. Restadh
\cite[Conjecture~1, p.~11]{diam} asks whether, whenever this operation
changes the imset, the segment joining the old and new imsets is
an edge of \(\CIM_G\).

\begin{theorem}
There is a five-vertex counterexample to the sink conjecture.
There is a connected counterexample of every order at least five.
An exhaustive exact check finds an exposing functional for every
imset-changing sink move on at most four vertices.
\end{theorem}

\section{Four orientations and their midpoint}
Take vertices \(0,1,2,3,4\) and the skeleton with edges
\[
  01,\quad 02,\quad 03,\quad 04,\quad 12,\quad 13.
\]
This is a diamond on \(0,1,2,3\) with a leaf \(4\) at \(0\).
Orient from earlier to later in each listed order. The entries of
the table are parent sets, with the dash denoting the empty set.
\[
\begin{array}{c|c|ccccc}
 &\text{order}&0&1&2&3&4\\ \hline
D&(4,2,0,3,1)&24&023&-&0&-\\
T&(4,0,3,1,2)&4&03&01&0&-\\
U&(4,0,3,2,1)&4&023&0&0&-\\
V&(4,2,0,1,3)&24&02&-&01&-
\end{array}
\]
Moving vertex \(2\) to the end changes precisely its two arcs,
so \(T=D\downarrow 2\). The pair coordinates are the same for all
four imsets. Among larger sets, precisely \(012\) and \(013\)
are nonzero in all four. The other nonzero coordinates are
\[
\begin{array}{c|l}
D&123,\ 0123,\ 024\\
T&\text{none}\\
U&123,\ 0123\\
V&024.
\end{array}
\]
Indeed, a parent set supplies the center together with each nonempty
subset of its parents. For example, vertex \(1\) of \(D\) contributes
\(012,013,123,0123\) in the displayed sizes. Vertex \(0\) of \(D\)
contributes \(024\).

The four imsets are distinct and satisfy
\[
                 c_D+c_T=c_U+c_V.
\]
If \([c_D,c_T]\) were an edge, its midpoint would be in that face.
The face property would put both \(c_U\) and \(c_V\) on the same
segment. The only binary points of a segment between different
binary vectors are its endpoints. This contradicts distinctness.
This use of a second pair with the same midpoint is the standard
rhombus criterion \cite[Sec.~2]{rhombus}.

\section{Connected lifts}
Take any graph containing the given five vertices as an induced
subgraph and with no new neighbor of \(2\). Place all new vertices
after the old ones in the four orders, using one common order
among the new vertices. Then each old parent set is unchanged,
and each new parent set is independent of which old order was used.
For a coordinate \(S\) containing a new vertex, a center in \(S\)
must be new, since no old vertex has a new parent. Thus every
coordinate involving a new vertex is common to all four imsets.
The old coordinates, midpoint equality, and distinctness are
preserved. So is the sink move, because \(2\) has no new incident
edge. Attaching the new vertices as leaves at \(4\) gives connected
examples of every larger order.

\section{Exact minimality check}
For \(1\le n\le4\), enumerate every labelled skeleton by choosing
its edge subset. Enumerate all acyclic orientations of a fixed
skeleton by trying each vertex order, since every DAG has a
topological order. Orienting from earlier to later and moving one
vertex to the end implements the sink operation exactly.

For each distinct pair of affected imsets \(a,b\) of a fixed
skeleton, let \(I\) be the coordinates where they differ.
Intersect the polytope with the constraints fixing all coordinates
outside \(I\) to their common binary value. This defines a face
because each condition is an equality in a valid bound \(0\le
x_S\le1\). The remaining vertices are exactly the distinct imsets
agreeing with these common values. On this face, seek integer weights
in \(\{-3,-2,\ldots,3\}^{I}\) that give equal score to \(a,b\) and
strictly smaller scores to every other remaining vertex. Existence
makes \([a,b]\) an edge of the face, hence of \(\CIM_G\).
The search checks actual integer inequalities.

For \(n=1,2\) no move changes an imset. At \(n=3,4\) there are
24 and 2208 changing moves before removing duplicates from orders,
direction and repeated choices. The corresponding counts of
distinct unordered imset pairs are \(3\) and \(202\). The replay
finds a supporting functional for each. Of the \(202\) pairs, 156
have no other vertices in their common-coordinate face. The
remaining 46 are certified with weights of magnitude at most 3.
This finite step is an independent exhaustive minimality check.

\clearpage
\section{Scope and provenance}
Linusson and Restadh \cite[Secs.~3 and 5]{rhombus} report the
rhombus criterion complete for \(\CIM_n\) through order five
and the full edge computation of \(\CIM_5\). That computation
necessarily already classified the five-vertex pair. The explicit
sink witness, its lift, and the direct small-order certificate are
the focus here. The tree case of the edge structure is known
\cite{edges}. The example has two triangles and reverses two arcs.
The result rules out the proposed sink-edge argument for a linear
diameter bound. The diameter bound itself remains a separate question.

\section*{Verification and technical disclosure}
The exact Python replay checks the imsets, nine larger connected
examples and minimality. A partial Lean~4 certificate checks the
five-vertex sink rewrite, distinct supports, coordinatewise
midpoint, and affected coordinates. The geometric argument and the
general lift are given here in ordinary mathematical proof.

Internally reviewed, not peer reviewed. Under human direction,
ChatGPT Astra~6.2 (OpenAI) generated the discovery proof and research
packet. OpenAI Codex independently checked the mathematics.
GPT-6.2 Sol (OpenAI) checked the literature and prepared the
manuscript, direct exposing-vector check, partial Lean certificate
and illustrations. DannyExperiments maintains the repository and is
the preferred citation identity. No AI system is proposed as an author.

\begin{thebibliography}{9}
\raggedright
\bibitem{diam} P. Restadh, \emph{Diameters of the characteristic imset
polytopes}, arXiv:2302.03647v2, 2023.
\bibitem{rhombus} S. Linusson and P. Restadh,
\emph{Rhombus criterion and the chordal graph polytope},
arXiv:2305.05275, 2023.
\bibitem{edges} S. Linusson, P. Restadh and L. Solus,
\emph{On the edges of characteristic imset polytopes},
Int. J. Approx. Reason. 191 (2026), 109606.
\end{thebibliography}
\end{document}
